% Options for packages loaded elsewhere
\PassOptionsToPackage{unicode}{hyperref}
\PassOptionsToPackage{hyphens}{url}
\PassOptionsToPackage{dvipsnames,svgnames,x11names}{xcolor}
%
\documentclass[
  letterpaper,
  DIV=11,
  numbers=noendperiod]{scrartcl}

\usepackage{amsmath,amssymb}
\usepackage{iftex}
\ifPDFTeX
  \usepackage[T1]{fontenc}
  \usepackage[utf8]{inputenc}
  \usepackage{textcomp} % provide euro and other symbols
\else % if luatex or xetex
  \usepackage{unicode-math}
  \defaultfontfeatures{Scale=MatchLowercase}
  \defaultfontfeatures[\rmfamily]{Ligatures=TeX,Scale=1}
\fi
\usepackage{lmodern}
\ifPDFTeX\else  
    % xetex/luatex font selection
\fi
% Use upquote if available, for straight quotes in verbatim environments
\IfFileExists{upquote.sty}{\usepackage{upquote}}{}
\IfFileExists{microtype.sty}{% use microtype if available
  \usepackage[]{microtype}
  \UseMicrotypeSet[protrusion]{basicmath} % disable protrusion for tt fonts
}{}
\makeatletter
\@ifundefined{KOMAClassName}{% if non-KOMA class
  \IfFileExists{parskip.sty}{%
    \usepackage{parskip}
  }{% else
    \setlength{\parindent}{0pt}
    \setlength{\parskip}{6pt plus 2pt minus 1pt}}
}{% if KOMA class
  \KOMAoptions{parskip=half}}
\makeatother
\usepackage{xcolor}
\setlength{\emergencystretch}{3em} % prevent overfull lines
\setcounter{secnumdepth}{5}
% Make \paragraph and \subparagraph free-standing
\makeatletter
\ifx\paragraph\undefined\else
  \let\oldparagraph\paragraph
  \renewcommand{\paragraph}{
    \@ifstar
      \xxxParagraphStar
      \xxxParagraphNoStar
  }
  \newcommand{\xxxParagraphStar}[1]{\oldparagraph*{#1}\mbox{}}
  \newcommand{\xxxParagraphNoStar}[1]{\oldparagraph{#1}\mbox{}}
\fi
\ifx\subparagraph\undefined\else
  \let\oldsubparagraph\subparagraph
  \renewcommand{\subparagraph}{
    \@ifstar
      \xxxSubParagraphStar
      \xxxSubParagraphNoStar
  }
  \newcommand{\xxxSubParagraphStar}[1]{\oldsubparagraph*{#1}\mbox{}}
  \newcommand{\xxxSubParagraphNoStar}[1]{\oldsubparagraph{#1}\mbox{}}
\fi
\makeatother


\providecommand{\tightlist}{%
  \setlength{\itemsep}{0pt}\setlength{\parskip}{0pt}}\usepackage{longtable,booktabs,array}
\usepackage{calc} % for calculating minipage widths
% Correct order of tables after \paragraph or \subparagraph
\usepackage{etoolbox}
\makeatletter
\patchcmd\longtable{\par}{\if@noskipsec\mbox{}\fi\par}{}{}
\makeatother
% Allow footnotes in longtable head/foot
\IfFileExists{footnotehyper.sty}{\usepackage{footnotehyper}}{\usepackage{footnote}}
\makesavenoteenv{longtable}
\usepackage{graphicx}
\makeatletter
\def\maxwidth{\ifdim\Gin@nat@width>\linewidth\linewidth\else\Gin@nat@width\fi}
\def\maxheight{\ifdim\Gin@nat@height>\textheight\textheight\else\Gin@nat@height\fi}
\makeatother
% Scale images if necessary, so that they will not overflow the page
% margins by default, and it is still possible to overwrite the defaults
% using explicit options in \includegraphics[width, height, ...]{}
\setkeys{Gin}{width=\maxwidth,height=\maxheight,keepaspectratio}
% Set default figure placement to htbp
\makeatletter
\def\fps@figure{htbp}
\makeatother

\KOMAoption{captions}{tableheading}
\makeatletter
\@ifpackageloaded{tcolorbox}{}{\usepackage[skins,breakable]{tcolorbox}}
\@ifpackageloaded{fontawesome5}{}{\usepackage{fontawesome5}}
\definecolor{quarto-callout-color}{HTML}{909090}
\definecolor{quarto-callout-note-color}{HTML}{0758E5}
\definecolor{quarto-callout-important-color}{HTML}{CC1914}
\definecolor{quarto-callout-warning-color}{HTML}{EB9113}
\definecolor{quarto-callout-tip-color}{HTML}{00A047}
\definecolor{quarto-callout-caution-color}{HTML}{FC5300}
\definecolor{quarto-callout-color-frame}{HTML}{acacac}
\definecolor{quarto-callout-note-color-frame}{HTML}{4582ec}
\definecolor{quarto-callout-important-color-frame}{HTML}{d9534f}
\definecolor{quarto-callout-warning-color-frame}{HTML}{f0ad4e}
\definecolor{quarto-callout-tip-color-frame}{HTML}{02b875}
\definecolor{quarto-callout-caution-color-frame}{HTML}{fd7e14}
\makeatother
\makeatletter
\@ifpackageloaded{caption}{}{\usepackage{caption}}
\AtBeginDocument{%
\ifdefined\contentsname
  \renewcommand*\contentsname{Table of contents}
\else
  \newcommand\contentsname{Table of contents}
\fi
\ifdefined\listfigurename
  \renewcommand*\listfigurename{List of Figures}
\else
  \newcommand\listfigurename{List of Figures}
\fi
\ifdefined\listtablename
  \renewcommand*\listtablename{List of Tables}
\else
  \newcommand\listtablename{List of Tables}
\fi
\ifdefined\figurename
  \renewcommand*\figurename{Figure}
\else
  \newcommand\figurename{Figure}
\fi
\ifdefined\tablename
  \renewcommand*\tablename{Table}
\else
  \newcommand\tablename{Table}
\fi
}
\@ifpackageloaded{float}{}{\usepackage{float}}
\floatstyle{ruled}
\@ifundefined{c@chapter}{\newfloat{codelisting}{h}{lop}}{\newfloat{codelisting}{h}{lop}[chapter]}
\floatname{codelisting}{Listing}
\newcommand*\listoflistings{\listof{codelisting}{List of Listings}}
\makeatother
\makeatletter
\makeatother
\makeatletter
\@ifpackageloaded{caption}{}{\usepackage{caption}}
\@ifpackageloaded{subcaption}{}{\usepackage{subcaption}}
\makeatother

\ifLuaTeX
  \usepackage{selnolig}  % disable illegal ligatures
\fi
\usepackage{bookmark}

\IfFileExists{xurl.sty}{\usepackage{xurl}}{} % add URL line breaks if available
\urlstyle{same} % disable monospaced font for URLs
\hypersetup{
  pdftitle={Git daily routine},
  pdfauthor={May Liu, Guillermo Peretz-Algorta},
  colorlinks=true,
  linkcolor={blue},
  filecolor={Maroon},
  citecolor={Blue},
  urlcolor={Blue},
  pdfcreator={LaTeX via pandoc}}


\title{Git daily routine}
\usepackage{etoolbox}
\makeatletter
\providecommand{\subtitle}[1]{% add subtitle to \maketitle
  \apptocmd{\@title}{\par {\large #1 \par}}{}{}
}
\makeatother
\subtitle{A walkthrough of Git practices}
\author{May Liu, Guillermo Peretz-Algorta}
\date{04-16-2025}

\begin{document}
\maketitle

\renewcommand*\contentsname{Table of contents}
{
\hypersetup{linkcolor=}
\setcounter{tocdepth}{3}
\tableofcontents
}

This document lists the daily routine steps of using Git. You don't have
to follow it exactly, but use it as a reference of what to keep in mind
when collaborating on projects with Git and GitHub.

\begin{tcolorbox}[enhanced jigsaw, toptitle=1mm, bottomtitle=1mm, arc=.35mm, titlerule=0mm, colframe=quarto-callout-tip-color-frame, colbacktitle=quarto-callout-tip-color!10!white, coltitle=black, bottomrule=.15mm, leftrule=.75mm, breakable, opacitybacktitle=0.6, title=\textcolor{quarto-callout-tip-color}{\faLightbulb}\hspace{0.5em}{Reminder}, colback=white, left=2mm, toprule=.15mm, opacityback=0, rightrule=.15mm]

In code that marks ``your\_branch'', replace it with the name of the
branch you're actually working on.

\end{tcolorbox}

\subsection{1. Fetch Latest Changes from
Remote}\label{fetch-latest-changes-from-remote}

Before you start working, pull in the latest remote changes without
modifying your working files.

\texttt{git\ fetch\ origin}

\textbf{What It Does:}

Fetching downloads new data (commits, branches, etc.) from the remote
repository, updating your local tracking branches. It does not alter
your working branches.

\textbf{Why It's Important:}

Keeping your tracking branches up to date helps you see if there are any
new changes (commits, merges, etc.) that might affect your work before
you merge.

\subsection{2. Check Which Branch You're
On}\label{check-which-branch-youre-on}

Ensure you're operating in the correct context by confirming your
current branch.

\texttt{git\ branch}

\textbf{What It Does:}

Lists all local branches with an asterisk indicating the active branch.

\textbf{Why It's Important:}

It's easy to lose track of what branch you're working on.
Double-checking prevents confusion and helps avoid making commits on the
wrong branch.

\subsection{3. Compare Your Current Branch with its Remote
Counterpart}\label{compare-your-current-branch-with-its-remote-counterpart}

A comparison reveals differences between what's in your local branch and
what's been pushed to the remote.

\texttt{git\ log\ HEAD..origin/your\_branch\ -\/-oneline}

\texttt{git\ log\ origin/your\_branch..HEAD\ -\/-oneline}

\textbf{What It Does:}

The first command shows commits present in the remote branch that
haven't been applied to your local branch.

The second command shows your local commits that haven't yet been pushed
to the remote branch.

\textbf{Why It's Important:}

This helps in catching up to remote changes or verifying your work has
been successfully pushed. It's a safeguard against merge conflicts and
duplicate work.

\subsection{4. Update Local Master with Latest Remote
Changes}\label{update-local-master-with-latest-remote-changes}

Get the latest `stable' branch changes from the remote before
integrating them into your feature branch.

\texttt{git\ switch\ master}

\texttt{git\ pull\ origin\ master}

\textbf{What It Does:}

Switches your working branch to master (or main).

Pulls and merges the latest changes from the remote master branch into
your local master branch.

\textbf{Why It's Important:}

Ensures that your base branch is current. A regularly updated base
minimizes conflicts when merging your feature branch later.

\subsection{5. Switch to Your Working
Branch}\label{switch-to-your-working-branch}

Go back to your feature or working branch where your work is in
progress.

\texttt{git\ switch\ your\_branch}

\textbf{What It Does:}

Changes your working directory to the specified branch.

\textbf{Why It's Important:}

It ensures that the next commands (e.g., merging) are performed on your
feature branch instead of master.

\subsection{6. Merge Updated Master into Your Working
Branch}\label{merge-updated-master-into-your-working-branch}

Incorporate the latest updates from master into your current work,
ensuring compatibility with the current project base.

\texttt{git\ merge\ master}

\textbf{What It Does:}

Merges the master branch changes into your current branch.

\textbf{Why It's Important:}

Integrating frequent updates from master can reduce large, complex merge
conflicts later, keeping your branch aligned with the main line of
development.

\subsection{7. Resolve Conflicts if Any}\label{resolve-conflicts-if-any}

If merge conflicts occur, you'll need to resolve them manually.

\textbf{What It Involves:}

Editing conflicting files to reconcile differences.

Testing to make sure that the resolutions do not break functionality.

\textbf{Why It's Important}

Proper resolution ensures stability and correctness of your codebase.

\subsection{8. Stage and Commit Your
Work}\label{stage-and-commit-your-work}

Once your changes are ready and conflicts are resolved, save them in
your local repository.

\texttt{git\ add\ .}

\texttt{git\ commit\ -m\ "Meaningful\ commit\ message"}

\textbf{What It Does:}

\texttt{git\ add\ .} stages all modified files.

\texttt{git\ commit} saves the changes with a descriptive message.

\textbf{Why It's Important:}

Structuring commits with clear messages improves project history, easing
code reviews and troubleshooting.

\subsection{9. Push Your Feature Branch to
Remote}\label{push-your-feature-branch-to-remote}

Publish your work online so it can be reviewed by your team members.

\texttt{git\ push\ origin\ your\_branch}

\textbf{What It Does:}

Pushes your committed changes to the remote repository (in the ``gpa''
branch).

\textbf{Why It's Important:}

This step makes your latest changes available to collaborators and
integrates your work into the shared repository.

\subsection{10. (Optional) Open a pull request to merge
branches}\label{optional-open-a-pull-request-to-merge-branches}

Coordinate with your team by formally requesting to merge your changes
into the master branch via a pull request (PR).

\textbf{What It Involves:}

Creating a PR on the repository hosting platform (like GitHub, GitLab,
or Bitbucket).

Engaging in code review and discussion about your changes.

\textbf{Why It's Important:}

Pull requests facilitate collaborative code review, ensure adherence to
coding standards, and help maintain code quality. :::




\end{document}
